\begin{abstract}
Effective waste management depends on correctly identifying an item's material, condition, disposal route, and handling threat. Existing computer-vision approaches largely emphasize material classification, leaving safe disposal under real-world conditions less studied. To address this gap, this thesis proposes a threat-aware, two-stage waste detection and classification system. A single-class YOLO11s detector localizes the item; a shared ConvNeXt-S backbone then predicts material (ten classes), condition (Good/Poor), route (Reuse, Recycle, Landfill, Special), and ordinal threat (0--3) using CORN. Calibration and regulation-grounded routing are evaluated alongside classification, with rules that allow predicted hazards to trigger special handling. Ten classifier configurations spanning convolutional networks, transformers, and an ensemble were compared, with ConvNeXt-S strongest among the individual backbones on material classification. Cross-validated evaluation over 457 items yields material accuracy 0.832 (macro-F1 0.803), threat QWK 0.771, condition AUROC 0.821, route macro-F1 0.675, and material ECE 0.036. The detector reaches mAP@50 0.992. On 50 test items, oracle and detector crops both yield material accuracy 0.840 at a detection rate of 1.0, with small decreases in condition and route performance. The routing policy sends 148 of 169 hazardous items to Special (recall 0.876); including two referrals for review raises safety recall to 0.888 at coverage 0.961. External validation and rare-class analysis identify remaining visual ambiguities, routing errors, and transfer limitations.

\textbf{Keywords:} Waste Classification, Object Detection, Multi-Task Learning, Ordinal Regression, Threat-Awareness, Calibration, YOLO, ConvNeXt, Computer Vision.
\end{abstract}
